\documentclass[11pt]{article}
\usepackage[margin=0.8in]{geometry}
\usepackage{amsmath,amssymb,booktabs}
\title{Separate genus-2 charts: checked generators and open descent}
\author{Formal result packet for independent review}
\date{25 October 2026}
\begin{document}
\maketitle
Let $A=K[X,Y]/(Y^2-X^5-a_3X^3-a_2X^2-a_1X-a_0)$ and let $B$ be the separately
presented infinity-chart coordinate ring. The candidate principal opens are
$A_X$ and $B_z$. The pinned Lean source defines a map from the \emph{polynomial
ring} $K[X,Y]$ to $B_z$ by $X\mapsto z^{-1}$ and $Y\mapsto wz^{-3}$.
Lean checks those generator images, $z^{-1}z=1$, the coefficient expansion
and the infinity-chart equation in $B_z$. Its compile exits zero, and the five
theorem axiom audits list only \texttt{propext}, \texttt{Classical.choice},
and \texttt{Quot.sound}.

\begin{center}\begin{tabular}{p{0.59\linewidth}p{0.29\linewidth}}\toprule
Gate & Status\\\midrule
Separate coordinate rings and principal opens & Defined\\
Polynomial map, generators, inverse, chart relation & Checked\\
Affine defining equation in the map kernel & Open\\
Descent to affine quotient and localization & Open\\
Reverse map and separate-chart ring isomorphism & Open\\
Scheme gluing, infinity stalk, divisor/$\operatorname{Pic}^0$ & Open\\
Rust law semantics and rho cost & Open / not measured\\\bottomrule\end{tabular}\end{center}

\noindent The precise next algebraic step is to derive the mapped affine
equation from the localized infinity relation and $zz^{-1}=1$, then use
\texttt{Ideal.Quotient.lift} and \texttt{IsLocalization.Away.lift}. A reverse
map and two composition proofs are still required. The prior same-ambient
quotient theorem does not identify these separately presented rings. See
\texttt{Common22SeparateChartMap.lean}, \texttt{separate-chart-map.csv}, and
\texttt{separate-chart-map.svg}. No empirical cost follows.
\end{document}
